\appendix


\chapter{Gamma ray flux from dark matter annihilation} \label{app:gammaflux}

In this appendix we derive the $\gamma$-ray flux from DM annihilation and the
associated $J$-factor introduced in \cref{Gammarays}.

The gamma ray flux produced by DM annihilation can be derived starting from the local DM density distribution.
The number of DM pairs contained in a volume $d^3r$ at position $\vec{r}$ is
\begin{equation} \label{dN}
    dN = d^3r \frac{1}{2} \left( \frac{\rho_{\mathrm{DM}}(\vec{r})}{m_{\mathrm{DM}}} \right)^2,
\end{equation}
where $\rho_{\mathrm{DM}}(\vec{r})$ is the DM mass density. We assume self-conjugate DM, i.e.\ the DM particles coincide with their antiparticles. On the other hand, in the case where DM particles can interact only with their own antiparticles, the factor $1/2$ must be replaced by $1/4$.
There is a quadratic dependence on the density since two DM particles are required for annihilation, whereas in the case of decay the dependence would be linear.
We introduce the factor $1/2$ to avoid double counting of particle pairs.

Assuming that a pair of DM particles annihilates into a pair of SM particles $f$, with a cross section $\langle \sigma v \rangle_f$, then multiplying \cref{dN} by $\langle \sigma v \rangle_f$ and summing over all possible final states, we obtain the number of annihilations per second in the volume $d^3r$

\begin{equation}
   \frac{dN}{dt} = d^3r \frac{1}{2} \left( \frac{\rho_{\mathrm{DM}}(\vec{r})}{m_{\mathrm{DM}}} \right)^2 \sum_f \langle \sigma v \rangle_f.
\end{equation}

Now multiplying this rate by the number of photons produced per annihilation in the channel $f$ within an infinitesimal energy bin $dE$, i.e.\ $\frac{dN_{f\to\gamma}}{dE}dE$, we obtain the number of photons produced per unit time and energy

\begin{equation}
   \frac{dN_{\gamma}}{dtdE} = d^3r \frac{1}{2} \left( \frac{\rho_{\mathrm{DM}}(\vec{r})}{m_{\mathrm{DM}}} \right)^2 \sum_f \langle \sigma v \rangle_f \frac{dN_{f\to\gamma}}{dE}. 
\end{equation}

As we show in \cref{GammaRay}, only a fraction of the photons produced at a distance $s$ from the observer reaches the telescope. Since the emission at each point is isotropic, the photons spread over a sphere of area $4\pi s^2$, so a detector of surface $S_{\mathrm{tel}}$ collects the geometric fraction
\begin{equation} \label{fraction}
\frac{S_{\mathrm{tel}}}{4\pi s^2}
\end{equation}
of the total number of photons produced. The solid angle $\Delta\Omega$ and the energy bin $\Delta E$ define the angular and spectral acceptance over which the signal is collected, and do not enter as a suppression factor.

\begin{figure}[H]
\centering
\includegraphics[width=\textwidth]{GammaRay.png}
\caption{\emph{A telescope of surface $S_{\mathrm{tel}}$ at line-of-sight distance $s$ collects the fraction $\frac{S_{\mathrm{tel}}}{4\pi s^2}$ of the photons produced by the isotropic emission \cite{GouttenoireSala2019}.}}
\label{GammaRay}
\end{figure}

The photon flux, i.e.\ the number of photons per unit detector area, time and energy, is obtained by multiplying the emission rate by the fraction \cref{fraction} and dividing by the collecting area $S_{\mathrm{tel}}$, which therefore cancels

\begin{equation}
\frac{d\Phi_\gamma}{dE} = \int \frac{1}{4\pi s^2}\, d^3r\; \frac{1}{2} \left( \frac{\rho_{\mathrm{DM}}(\vec{r})}{m_{\mathrm{DM}}} \right)^2 \sum_f \langle \sigma v \rangle_f \frac{dN_{f\to\gamma}}{dE}.
\end{equation}

Writing the volume element seen by the observer as $d^3r = s^2\, ds\, d\Omega$, with $s$ the coordinate along the line of sight, and integrating over the line of sight we obtain the differential photon flux per unit solid angle and energy

\begin{equation} \label{Annihilation}
    \frac{d\Phi_\gamma}{d\Omega dE} (\theta) = \frac{1}{2} \frac{r_\odot}{4\pi} \left( \frac{\rho_\odot}{m_{\mathrm{DM}}} \right)^2 J_{\mathrm{ann}}(\theta) \sum_f \langle \sigma v \rangle_f \frac{dN_{f\to\gamma}}{dE},
\end{equation}
where in the particular case of the Milky Way

\begin{equation}
    J_{\mathrm{ann}} (\theta) = \int_{l.o.s} \frac{ds}{r_\odot} \left( \frac{\rho_{\mathrm{DM}}(r(s, \theta))}{\rho_\odot} \right)^2
\end{equation}
is the annihilation $J$-factor. 
Here $s$ is the distance from the observer along the line of sight, $\theta$ is the angle between the line of sight and the direction of the Galactic Center, and $r$ is the galactocentric radius, i.e.\ the distance from the Galactic Center that sets the DM density $\rho_{\mathrm{DM}}(r)$;
it is related to $s$ and $\theta$ by $r = \sqrt{r_\odot^{2} + s^{2} - 2\,r_\odot\, s \cos\theta}\,$. Finally, $r_\odot \simeq 8.33$~kpc is the distance from the Earth to the Galactic Center and $\rho_\odot \simeq 0.4~\text{GeV/cm}^3$ is the local DM mass density.

Note that a similar treatment for DM decay leads to

\begin{equation} \label{Decay}
\frac{d\Phi_\gamma}{d\Omega\, dE} (\theta)
=
\frac{r_\odot}{4\pi}
\left(\frac{\rho_\odot}{m_{\mathrm{DM}}}\right)
J_{\mathrm{dec}}(\theta)
\sum_{f}\Gamma_f\,
\frac{dN_{f\to\gamma}}{dE}\,,
\end{equation}
where

\begin{equation} 
J_{\mathrm{dec}} (\theta)
=
\int_{l.o.s} \frac{ds}{r_\odot}
\left(\frac{\rho_{\mathrm{DM}}(r(s, \theta))}{\rho_\odot}\right)\,
\label{eq:Jdec}
\end{equation}
is the $J$-factor for decay and $\Gamma_f$ is the decay rate in the channel $f$. As expected, we obtain a linear dependence on the DM density.










\chapter{Phase-space calculations for the adiabatic spike}
\label{app:spike}

In this appendix we report all the explicit derivations used in the study of the DM spike in \cref{chap:spikes}.

\section{The density as an integral over energy and angular momentum}
\label{app:densityEL}

Here we derive \cref{eq:density_phase_space} of \cref{adiabatic_spike}.

In a spherically symmetric system it is useful to decompose the
velocity into its radial $v_r$ and tangential $v_t$ components. 
In cylindrical coordinates, recalling once more the invariance of $f$ under rotations,
the volume element is $d^3v = 2\pi\, v_t\, dv_t\, dv_r$, where we already integrated over the azimuthal angle.
The density becomes

\begin{equation} \label{density_velocities}
\rho'(r) = \int f'\, d^3v = 2\pi \int f'\, v_t\, dv_t\, dv_r .
\end{equation}

We now change variables again from $(v_r, v_t)$ to $(E', L')$, defined by

\begin{equation}
E' = \phi'(r) + \frac{v_r^2 + v_t^2}{2},
\qquad
L' = r\, v_t ,
\end{equation}
where $\phi'$ is the gravitational potential and the primes indicate that these quantities are evaluated after the formation of the black hole.

At fixed $r$, the Jacobian is

\begin{equation}
\frac{\partial(E', L')}{\partial(v_r, v_t)}
=
\det
\begin{pmatrix}
v_r & v_t \\
0 & r
\end{pmatrix}
= r\, v_r
\qquad\Longrightarrow\qquad
dv_r\, dv_t = \frac{dE'\, dL'}{r\, v_r} .
\end{equation}

Now we notice that this change of coordinates is not injective: 
to a given $(E',L')$ there correspond two velocities: the
orbit crosses the radius $r$ once moving outward ($v_r > 0$) and once moving
inward ($v_r < 0$), with the same $|v_r|$.
Hence we have to account for both by multiplying the integral by a factor of two

\begin{equation}
\rho'(r)
=
2 \times 2\pi \int \frac{L'}{r}\, \frac{dE'\, dL'}{r\, v_r}\, f'(E',L')
=
\int dE' \int dL'\, \frac{4\pi L'}{r^2 v_r}\, f'(E',L') .
\end{equation}

The integration limits follow from the
requirements that the orbit is bound ($E' < 0$), that it is not captured by
the BH ($L' > L'_c$), and that $v_r^2 \ge 0$ ($L' < L'_m$).




\section{Potential of the power-law halo}
\label{app:powerlawpot}

This derivation completes the potential used in \cref{sec:generalscaling}.

For the density profile

\begin{equation}
    \rho(r) = \rho_s \left( \frac{r}{r_s} \right)^{-\gamma},
    \qquad 0 \le \gamma < 2,
\end{equation}
the enclosed mass is

\begin{equation} \label{eq:enclosedmass}
M_h(r) = \int_0^r 4\pi r'^2 \rho(r')\, dr'
= \frac{4\pi \rho_s r_s^{\gamma}}{3-\gamma}\, r^{3-\gamma},
\end{equation}
which converges at the center for $\gamma < 3$.

We now recall Poisson's equation for a spherically symmetric system

\begin{equation}
\nabla^2 \phi = \frac{1}{r^2}\frac{d}{dr}\!\left( r^2 \frac{d\phi}{dr} \right) 
= 4\pi G\, \rho(r).
\end{equation}

Integrating this equation we obtain Newton's theorem

\begin{equation}
r^2 \frac{d\phi}{dr} = \int_0^r 4\pi G\, \rho(r')\, r'^2\, dr' = G\, M_h(r)
\qquad\Longrightarrow\qquad
\frac{d\phi}{dr} = \frac{G\, M_h(r)}{r^2}.
\end{equation}

Now, substituting $M_h(r)$ from \cref{eq:enclosedmass} and integrating once more, 
with the boundary condition $\phi(0)=0$, we derive the potential


\begin{equation} \label{eq:app_potential}
\phi(r) = \frac{4\pi G \rho_s r_s^{\gamma}}{3-\gamma}\int_0^r r'^{\,1-\gamma}\, dr'
= \phi_0 \left( \frac{r}{r_s} \right)^{2-\gamma},
\qquad
\phi_0 = \frac{4\pi G \rho_s r_s^{2}}{(3-\gamma)(2-\gamma)} .
\end{equation}

Note that the last integral, $\int_0^r r'^{\,1-\gamma}\, dr' = r^{2-\gamma}/(2-\gamma)$, converges only for 
$\gamma < 2$. Therefore the condition at the origin $\phi(0)=0$ is admissible.

Another important consideration is that, although derived here for the single power law,
this result has a wider validity: for any spherically symmetric density profile with inner slope $\gamma$,
the potential has the same central behavior, determined by the enclosed mass.




\section{Eddington's inversion formula}
\label{app:eddington}

Here we derive Eddington's formula used in \cref{sec:generalscaling}.

We want now to obtain the phase-space distribution as a function of the density profile $\rho$,
the opposite of what we did in appendix~\ref{app:densityEL}.

We start again from the density as an integral over velocity space, 
$\rho = \int f\, d^3v$.
Now we have an isotropic distribution $f(E)$, i.e.\ $f$ depends only on the energy,
so the volume element in spherical coordinates reduces to $d^3v = 4\pi\, v^2\, dv$. 
Changing variables through $E = \phi(r) + v^2/2$, the density becomes

\begin{equation}
\rho = 4\pi \int_0^{\infty} f(E)\, v^2\, dv
= 4\pi\sqrt{2} \int_{\phi}^{\infty} f(E)\, \sqrt{E - \phi}\; dE ,
\label{eq:app_rhof}
\end{equation}
where we used $v = \sqrt{2(E-\phi)}$ and $v^2 dv = \sqrt{2}\sqrt{E-\phi}\, dE$.

Since both $\rho$ and $\phi$ are functions of $r$, and $\phi(r)$ is monotonic 
in the region of interest, we can regard
$\rho$ as a function of $\phi$ and differentiate both sides with respect to $\phi$

\begin{equation}
\frac{d\rho}{d\phi}
= -2\sqrt{2}\,\pi \int_{\phi}^{\infty} \frac{f(E)}{\sqrt{E-\phi}}\; dE .
\label{eq:app_abel}
\end{equation}

\Cref{eq:app_abel} is an Abel integral equation for $f$.

To invert it, we multiply by
$(\phi - \mathcal{E})^{-1/2}$, with $\mathcal{E}$ as a new fixed parameter, and integrate over $\phi$ from $\mathcal{E}$ to
infinity, exchanging the order of integration:
\begin{equation} \label{rho_eddington}
\int_{\mathcal{E}}^{\infty} \frac{d\rho}{d\phi}\,
\frac{d\phi}{\sqrt{\phi-\mathcal{E}}}
= -2\sqrt{2}\,\pi \int_{\mathcal{E}}^{\infty} dE\, f(E)
\int_{\mathcal{E}}^{E} \frac{d\phi}{\sqrt{(\phi-\mathcal{E})(E-\phi)}} .
\end{equation}

The inner integral

\begin{equation}
\int_{\mathcal{E}}^{E} \frac{d\phi}{\sqrt{(\phi-\mathcal{E})(E-\phi)}} 
\end{equation}
is a standard integral, independent of both $\mathcal{E}$ and $E$, which can be 
solved explicitly with the substitution $\phi = \mathcal{E} + (E-\mathcal{E})\sin^2\theta$
and becomes

\begin{equation}
    2\int_0^{\pi/2} d\theta = \pi.
\end{equation}

Now, substituting this value in \cref{rho_eddington}

\begin{equation}
\int_{\mathcal{E}}^{\infty} \frac{d\rho}{d\phi}\,
\frac{d\phi}{\sqrt{\phi-\mathcal{E}}}
= -2\sqrt{2}\,\pi^2 \int_{\mathcal{E}}^{\infty} f(E)\, dE ,
\end{equation}
and differentiating with respect to $\mathcal{E}$ we obtain Eddington's
formula 

\begin{equation}
f(E)
= \frac{1}{\sqrt{8}\,\pi^2}\,
\frac{d}{dE} \int_{E}^{\infty} \frac{d\rho}{d\phi}\,
\frac{d\phi}{\sqrt{\phi-E}} .
\label{eq:app_eddington}
\end{equation}



\section{Radial actions}
\label{app:actions}

Here we derive the radial actions used in \cref{sec:generalscaling}.

\subsection*{Power-law potential}

For the initial potential

\begin{equation}
    \phi(r) = \phi_0 \left( \frac{r}{r_s} \right)^{2-\gamma}
\end{equation}
the action integral cannot be performed in closed form for generic $\gamma$, 
so we will study its scaling with $E$.

We can obtain the orbit $r_E$ by $\phi(r_E) = E$:

\begin{equation}
    r_E = r_s \left( \frac{E}{\phi_0} \right)^{\frac{1}{2-\gamma}}.
\end{equation}

Recalling the radial action expression in \cref{radialaction}
and rescaling $r = r_E\, x$ we obtain

\begin{equation}
I(E,L)
= 2 \int_{r_-}^{r_+} \sqrt{2E - 2\phi(r) - \frac{L^2}{r^2}}\; dr
= 2\sqrt{2E}\; r_E \int_{x_-}^{x_+}
\sqrt{1 - x^{2-\gamma} - \frac{\lambda^2}{x^2}}\; dx ,
\end{equation}
where $\lambda^2 = \frac{L^2}{2E\, r_E^2}$ is dimensionless.

Hence

\begin{equation}
I(E,L) = 2\sqrt{2E}\; r_E \; F(\lambda)
\;\propto\; E^{\frac{1}{2} + \frac{1}{2-\gamma}}
= E^{\frac{4-\gamma}{2(2-\gamma)}}
\end{equation}
at fixed $\lambda$. 




\subsection*{Point-mass potential}

The final potential is the point-mass

\begin{equation}
    \phi'(r) = -\frac{GM_{\mathrm{BH}}}{r},
\end{equation}
so the radial velocity becomes

\begin{equation}
v_r^2 = 2E' + \frac{2GM_{\mathrm{BH}}}{r} - \frac{L'^2}{r^2}
= \frac{2|E'|}{r^2} \left( r_+ - r \right)\left( r - r_- \right),
\end{equation}
where $r_\pm$ are the turning points, with

\begin{equation}
r_+ + r_- = \frac{GM_{\mathrm{BH}}}{|E'|},
\qquad
r_+ r_- = \frac{L'^2}{2|E'|} .
\end{equation}

The radial action is therefore

\begin{equation}
I' = 2 \int_{r_-}^{r_+} v_r\, dr
= 2\sqrt{2|E'|} \int_{r_-}^{r_+}
\frac{\sqrt{(r_+ - r)(r - r_-)}}{r}\; dr .
\end{equation}

The definite integral is a standard one and can be evaluated with the
substitution $r = r_+ \sin^2\theta + r_- \cos^2\theta$:

\begin{equation}
\int_{r_-}^{r_+} \frac{\sqrt{(r_+-r)(r-r_-)}}{r}\, dr
= \frac{\pi}{2}\left( r_+ + r_- \right) - \pi \sqrt{r_+ r_-} ,
\end{equation}
so that

\begin{equation}
I' = 2\sqrt{2|E'|} \left[ \frac{\pi}{2}\frac{GM_{\mathrm{BH}}}{|E'|}
- \pi \frac{L'}{\sqrt{2|E'|}} \right]
= 2\pi \left( -L' + \frac{GM_{\mathrm{BH}}}{\sqrt{-2E'}} \right) .
\end{equation}

Note that the action vanishes for circular orbits, for which
$L' = \frac{GM_{\mathrm{BH}}}{\sqrt{-2E'}}$ is the maximum angular momentum at fixed energy.




\section{Velocity integral for the spike density}
\label{app:spikevelocity}

Here we perform the velocity integral quoted in \cref{sec:generalscaling}.

We saw that in the region dominated by the black hole the potential is

\begin{equation}
    \phi'(r) \simeq -GM_{\mathrm{BH}}/r,
\end{equation}
and bound particles have 

\begin{equation}
E' = -\frac{GM_{\mathrm{BH}}}{r} + \frac{v^2}{2},
\qquad
0 \le v \le v_{\mathrm{esc}} = \sqrt{\frac{2GM_{\mathrm{BH}}}{r}} .
\end{equation}

The final distribution function, obtained from adiabatic invariance considerations, scales as

\begin{equation}
    f'(E') \propto (-E')^{\,n\frac{2-\gamma}{4-\gamma}}.
\end{equation}

So, the density is given by the integral over velocity space

\begin{equation}
\rho'(r)
= 4\pi \int_0^{v_{\mathrm{esc}}} f'(E')\, v^2\, dv
\;\propto\;
\int_0^{v_{\mathrm{esc}}}
\left( \frac{GM_{\mathrm{BH}}}{r} - \frac{v^2}{2} \right)^{\! n\frac{2-\gamma}{4-\gamma}} v^2\, dv .
\end{equation}


With the substitution $v = v_{\mathrm{esc}}\, u$, $u \in [0,1]$, we can
factor out all the dependence on $r$

\begin{equation}
\rho'(r)
\;\propto\;
\left( \frac{GM_{\mathrm{BH}}}{r} \right)^{n\frac{2-\gamma}{4-\gamma}}
v_{\mathrm{esc}}^{3}
\int_0^1 \left(1-u^2\right)^{n\frac{2-\gamma}{4-\gamma}} u^2\, du
\;\propto\;
\left( \frac{GM_{\mathrm{BH}}}{r} \right)^{n\frac{2-\gamma}{4-\gamma}+\frac{3}{2}}
\;\propto\;
r^{-\left(\frac{3}{2} + n\frac{2-\gamma}{4-\gamma}\right)} ,
\end{equation}
where the integral in $u$ is a finite constant independent of $r$, in particular we notice that

\begin{equation}
\int_0^1 \left(1-u^2\right)^{n\frac{2-\gamma}{4-\gamma}} u^2\, du
= \frac{1}{2}\int_0^1 t^{1/2}\,(1-t)^{n\frac{2-\gamma}{4-\gamma}}\, dt
= \frac{1}{2}\,
B\!\left(\frac{3}{2},\; n\,\frac{2-\gamma}{4-\gamma}+1\right),
\end{equation}
where $B$ is Euler's Beta function, whose definition will be better discussed in appendix~\ref{app:fGS}.

We therefore obtain
\begin{equation}
\rho'(r) \propto r^{-\gamma_{\mathrm{sp}}},
\qquad
\gamma_{\mathrm{sp}} = \frac{3}{2} + n\,\frac{2-\gamma}{4-\gamma} .
\end{equation}






\section{Distribution function of the power-law halo}
\label{app:fGS}

Here we derive the Gondolo--Silk distribution function used in \cref{sec:GSspike}.

For the power law halo

\begin{equation} \label{eq:app_exactpowerlaw}
   \rho(r) = \rho_s \left( \frac{r}{r_s} \right)^{-\gamma},
    \qquad 0<\gamma<2,
\end{equation}
it is possible to evaluate Eddington's formula exactly.

We recall Eddington's formula of \cref{eq:app_eddington}

\begin{equation} \label{eq:app_exacteddington}
f(E)
= \frac{1}{\sqrt{8}\,\pi^2}\,
\frac{d}{dE} \int_{E}^{\infty} \frac{d\rho}{d\phi}\,
\frac{d\phi}{\sqrt{\phi-E}},
\end{equation}
and the potential of \cref{eq:app_potential}, which is exact for the considered density profile,

\begin{equation} \label{eq:app_exactpotential}
\phi(r) = \frac{4\pi G \rho_s r_s^{\gamma}}{3-\gamma}\int_0^r r'^{\,1-\gamma}\, dr'
= \phi_0 \left( \frac{r}{r_s} \right)^{2-\gamma},
\qquad
\phi_0 = \frac{4\pi G \rho_s r_s^{2}}{(3-\gamma)(2-\gamma)} .
\end{equation}

Combining \cref{eq:app_exactpowerlaw} with \cref{eq:app_exactpotential} we can eliminate the radius

\begin{equation}
\rho(\phi) = \rho_s \left( \frac{\phi}{\phi_0} \right)^{-p},
\qquad
p \equiv \frac{\gamma}{2-\gamma},
\qquad
\frac{d\rho}{d\phi} = -\frac{p\,\rho_s}{\phi_0}
\left( \frac{\phi}{\phi_0} \right)^{-p-1} .
\end{equation}

Now the integral in \cref{eq:app_exacteddington} can be computed with the
substitution $\phi = E/x$, $x \in (0,1]$:

\begin{equation} \label{beta}
\int_{E}^{\infty} \phi^{-p-1}\, (\phi-E)^{-1/2}\, d\phi
= E^{-p-1/2} \int_0^1 x^{\,p-1/2} (1-x)^{-1/2}\, dx
= E^{-p-1/2}\, B\!\left(p+\tfrac12, \tfrac12\right),
\end{equation}
where $B$ is Euler's Beta function, defined for $a,b > 0$ as

\begin{equation}
B(a,b) = \int_0^1 x^{\,a-1}(1-x)^{\,b-1}\, dx
= \frac{\Gamma(a)\,\Gamma(b)}{\Gamma(a+b)},
\end{equation}
with $\Gamma$ Euler's Gamma function,

\begin{equation} \label{eq:app_gammadef}
\Gamma(z) = \int_0^{\infty} t^{\,z-1} e^{-t}\, dt ,
\qquad z > 0 .
\end{equation}

The convergence of the integral in \cref{beta} is guaranteed
because $p + \tfrac12 > 0$, since $p = \frac{\gamma}{2-\gamma} > 0$ for $0 < \gamma < 2$.

Taking the derivative with respect to $E$, we obtain

\begin{equation}
f(E)
= \frac{\rho_s\, \phi_0^{\,p}}{\sqrt{8}\,\pi^2}\;
p \left(p+\tfrac12\right) B\!\left(p+\tfrac12,\tfrac12\right)
E^{-p-3/2} ,
\end{equation}
so that $f(E) \propto E^{-n}$ with
\begin{equation}
n = p + \frac{3}{2} = \frac{6-\gamma}{2(2-\gamma)} .
\end{equation}

The coefficient can be simplified using the properties of Euler's Beta and Gamma functions:

\begin{equation} \label{eq:app_gammaprop}
\Gamma(z+1) = z\,\Gamma(z),
\qquad
\Gamma\!\left(\tfrac12\right) = \sqrt{\pi} ,
\end{equation}
and in our particular case

\begin{equation}
B\!\left(p+\tfrac12,\tfrac12\right)
= \sqrt{\pi}\,\frac{\Gamma\!\left(p+\tfrac12\right)}{\Gamma(p+1)} ,
\end{equation}
together with the identities

\begin{equation}
\left(p+\tfrac12\right)\Gamma\!\left(p+\tfrac12\right)
= \Gamma\!\left(p+\tfrac32\right) = \Gamma(n),
\qquad
\frac{p}{\Gamma(p+1)} = \frac{1}{\Gamma(p)} = \frac{1}{\Gamma\!\left(n-\tfrac32\right)} .
\end{equation}

Finally we obtain the Gondolo--Silk distribution function

\begin{equation}
f(E)
= \frac{\rho_s}{(2\pi\phi_0)^{3/2}}\,
\frac{\Gamma(n)}{\Gamma\!\left(n-\tfrac32\right)}
\left( \frac{\phi_0}{E} \right)^{n} .
\end{equation}





\section{Scaling of the singular terms in Eddington's formula}
\label{app:nkscaling}

Here we derive the scaling used in \cref{sec:gamma0}.

We recall the expansion of Eddington's integral of \cref{eq:eddexp}


\begin{equation} \label{eq:app_eddexp}
f(E)
\sim
\int_{r_\star}^{R} \frac{dr}{r^{3}}\,
\frac{c_0 + c_1 (r/r_s)^{\alpha} + c_2 (r/r_s)^{2\alpha} + \dots}
{\sqrt{\, r^{2-\gamma} - E \,}},
\end{equation}
where the lower limit $r_\star = E^{1/(2-\gamma)}$ is the radius at which
$\phi(r_\star) = E$, while $R$ is an arbitrary upper radius, of the order
of the scale radius of the halo.
We set $\phi_0 = r_s = 1$, since we
are only interested in the scaling with $E$.


A term $\propto r^{k}$ in the numerator gives the contribution


\begin{equation} \label{eq:app_Ik}
I_k(E) = \int_{r_\star}^{R}
\frac{r^{\,k-3}}{\sqrt{r^{2-\gamma} - E}}\, dr .
\end{equation}

Rescaling $r = r_\star x$ and using $r_\star^{2-\gamma} = E$, so that
$r^{2-\gamma} - E = E\,(x^{2-\gamma}-1)$, we obtain

\begin{equation} \label{eq:app_nk}
I_k(E) = r_\star^{\,k-2}\, E^{-1/2}
\int_1^{R/r_\star} \frac{x^{\,k-3}\, dx}{\sqrt{x^{2-\gamma}-1}}
\;\propto\; E^{\frac{k-2}{2-\gamma} - \frac{1}{2}}
= E^{-n_k},
\qquad
n_k = \frac{6-\gamma-2k}{2(2-\gamma)} ,
\end{equation}
then the scaling is all in the prefactor, provided that the dimensionless integral in $x$ tends to a finite constant as
$E \to 0^{+}$. We therefore examine its two endpoints.

At the lower endpoint the integrand is singular, but integrably, so
expanding $x^{2-\gamma}-1 \simeq (2-\gamma)(x-1)$ for $x \to 1^{+}$, it
behaves as $(x-1)^{-1/2}$ and never produces a divergence of $f(E)$.

At the upper endpoint, where $E \to 0^{+}$ implies $R/r_\star \to \infty$,
the integrand behaves as $x^{\,k-3-(2-\gamma)/2}$, so the integral
converges if and only if

\begin{equation} \label{eq:app_conv}
k - 3 - \frac{2-\gamma}{2} < -1
\quad \Longleftrightarrow \quad
n_k > 0 ,
\end{equation}
i.e.\ the condition of convergence at large $x$ is precisely the condition
that the exponent $n_k$ be positive.
Instead, when $n_k \le 0$, the
integral is dominated by large radii and $I_k$
contributes only to the regular part
of $f(E)$, as discussed in \cref{sec:gamma0}.


For $k = 0$ one recovers $n_0 = \frac{6-\gamma}{2(2-\gamma)}$, in agreement with the exact result of
Gondolo and Silk.
For a cored profile, with $\gamma = 0$ and $k = \alpha$, one obtains $n_\alpha = \frac{3-\alpha}{2}$, 
which is the case analysed in Ref.~\cite{Ullio_2001}.
For the non-singular isothermal
profile, with $\gamma = 0$ and $\alpha = 2$, both $c_0$ and $c_1$ vanish
and the first surviving term has $k = 2\alpha = 4$, for which
$n_4 = -1/2 < 0$: no term of the expansion produces a divergence, hence
$f(E)$ is regular and $n = 0$.






\chapter{Saturation density from dark matter annihilation}
\label{app:saturation}

In this appendix we derive the saturation density of \cref{sec:saturation}.

The density
growth in the innermost region of the DM spike is depleted because DM annihilates efficiently
on account of very high DM density. The evolution of the density due to annihilation is

\begin{equation}
    \frac{\partial \rho}{\partial t} = -\frac{\langle\sigma v\rangle}{m_{\mathrm{DM}}}\,\rho^2,
    \label{eq:app_evolution}
\end{equation}
where $m_{\mathrm{DM}}$ is the DM mass and $\langle\sigma v\rangle$ is the cross section.

\Cref{eq:app_evolution} is a separable first-order differential equation, so writing it as

\begin{equation}
    \frac{d\rho}{\rho^2} = -\frac{\langle\sigma v\rangle}{m_{\mathrm{DM}}}\, dt,
\end{equation}
and integrating from the initial density $\rho_i$ at the time of
black hole formation up to the density $\rho$ at a later time $t$

\begin{equation}
    \int_{\rho_i}^{\rho} \frac{d\rho'}{\rho'^2}
    = -\frac{\langle\sigma v\rangle}{m_{\mathrm{DM}}}\int_0^t dt',
\end{equation}
we obtain 

\begin{equation}
    \frac{1}{\rho} = \frac{1}{\rho_i} + \frac{\langle\sigma v\rangle}{m_{\mathrm{DM}}}\, t.
    \label{eq:app_solution}
\end{equation}

Considering a very large initial density we can neglect the term $1/\rho_i$ and the equation becomes

\begin{equation} \label{DMdensity-time}
    \rho(t) \simeq \frac{m_{\mathrm{DM}}}{\langle\sigma v\rangle\, t}.
\end{equation}

Finally, we obtain the saturation density evaluating \cref{DMdensity-time} at the age of the black hole $t = t_{\mathrm{BH}}$, conservatively $10^{10}$ yr,
the maximum time over which annihilation has been acting,

\begin{equation}
    \rho_{\mathrm{core}} = \frac{m_{\mathrm{DM}}}{\langle\sigma v\rangle\, t_{\mathrm{BH}}}.
\end{equation}

$\rho_{\mathrm{core}}$ can be obtained also as the density at which the
characteristic annihilation timescale
$t_{\mathrm{ann}} \sim m_{\mathrm{DM}}/(\langle\sigma v\rangle\,\rho)$ becomes comparable to
$t_{\mathrm{BH}}$.
